% !TEX TS-program = xelatex
\documentclass[10.5pt,a4paper]{ctexart}
\usepackage[a4paper,margin=1.65cm]{geometry}
\usepackage{titlesec}
\usepackage{enumitem}
\usepackage{hyperref}
\usepackage{xcolor}
\usepackage{microtype}
\definecolor{accent}{HTML}{1F4E79}
\hypersetup{colorlinks=true,urlcolor=accent,linkcolor=accent}
\setCJKmainfont{Hiragino Sans GB}
\setmainfont{Times New Roman}
\setlength{\parindent}{0pt}
\setlength{\parskip}{2pt}
\setlist[itemize]{leftmargin=1.5em,itemsep=1pt,topsep=2pt,parsep=0pt}
\titleformat{\section}{\large\bfseries\color{accent}}{}{0pt}{}[\vspace{-4pt}\color{accent}\rule{\linewidth}{0.5pt}]
\titlespacing*{\section}{0pt}{8pt}{4pt}
\newcommand{\entry}[2]{\textbf{#1}\hfill #2\\}
\begin{document}
\begin{center}
{\LARGE\bfseries 肖童仁（Tongren Xiao）}\\[3pt]
中山大学数学学院 \quad $\cdot$ \quad 广州，中国\\
\href{mailto:tongrenxiao7@gmail.com}{tongrenxiao7@gmail.com} \quad $\cdot$ \quad
\href{https://xtr1976536.github.io/}{xtr1976536.github.io}
\end{center}

\section*{研究兴趣}
金融波动率建模与预测；金融时间序列的形状空间方法与基于形状的表示；预测校准与低估风险控制；Schubert 演算与代数几何；部分困难数学猜想的研究。

\section*{教育经历}
\entry{中山大学数学学院}{2023--2027年9月（预计）}
数学专业本科生\\
GPA：3.00/4.00；平均分：80/100

\section*{科研经历}
\entry{量化金融研究}{2025年至今}
指导教师：黄怡蓉副教授，中山大学\\
从事金融波动率建模与预测研究，涉及混合 LSTM--GARCH 比较、实现波动率的几何类比预测，以及针对预测低估风险的后处理校准方法。

\entry{Schubert 演算研究}{2025年至今}
指导教师：李长征教授，中山大学\\
研究量子双 Schubert 多项式的正性及 Schubert 演算中的相关结构。

\entry{对部分困难数学猜想的研究}{2026年至今}
围绕矩阵分析中的三框架子矩阵问题和数论中的 Erdős--Straus 猜想，研究固定平移障碍、大筛法估计、Dahan 筛法常数优化以及严格的部分结果和剩余间隙。

\entry{金融领域世界模型研究}{2026年至今}
\textit{Anchor-Conditioned Probabilistic World Models for Multi-Asset Realized-Volatility Paths}（进行中）。探索潜在状态空间动力学、HAR--IV 锚定、多步残差似然和金融波动率路径的联合现实性。

\section*{课程学习与学术讲授}
选修两门研究生课程、一门荣誉课程和一门跨专业课程：论文选读、群上的调和分析、代数拓扑、随机过程。在论文选读课程中负责讲授 Emily Riehl《Category Theory in Context》第 1 章和第 3 章。

\section*{国际化学术训练}
\entry{Introduction to Option Pricing}{2025年3月7日--5月25日}
授课教师：Johannes Ruf；成绩：81.50。

\section*{代表性论文与研究笔记}
\begin{itemize}
  \item \textbf{Shape Is Useful: Geometric Analog Forecasting of Realized Volatility.} Tongren Xiao and Shuhang Chen，手稿，2026。
  \item \textbf{A Comparative Study of Hybrid LSTM Frameworks for Volatility Forecasting in the NASDAQ-100 and S\&P 500 Markets.} 手稿，2025。
  \item \textbf{Graham Positivity of Quantum Double Schubert Polynomials.} Tongren Xiao and Yihua Yue，手稿，2026。
  \item \textbf{On the Submatrices with the Best-Bounded Inverses for Matrices with Three Orthonormal Columns.} 研究笔记，2026。
  \item \textbf{Fixed-Shift Obstructions and Optimized Sieve Constants for the Erdős--Straus Conjecture.} 研究笔记，2026。
\end{itemize}

\section*{荣誉与项目}
\begin{itemize}
  \item 2026 COMAP 数学建模竞赛（MCM）B 题 Honorable Mention。
  \item 2025 全国大学生数学建模竞赛广东赛区 A 题二等奖。
  \item 省级大学生创新创业训练计划，项目编号 20251647：\textit{基于多期权组合策略的资产波动率建模}。
\end{itemize}

\section*{技能}
熟悉 LaTeX 学术写作与论文排版；熟悉 MATLAB 数学建模与数值计算；数学证明与技术写作；线性代数、矩阵分析和抽象代数结构。
\end{document}
